\documentclass[aspectratio=169,10pt,spanish]{beamer}

\input{../../../_preambulo_beamer.tex}
\input{../../../_comandos_beamer.tex}

\selectlanguage{spanish}

\title{MA1001B - Modelación Estadística para la Toma de Decisiones}
\subtitle{Lección 40: Bondad de ajuste ji-cuadrada}
\author{}
\institute{Tecnol\'ogico de Monterrey}
\date{}

\begin{document}

\begin{frame}[plain]
  \vspace{1.2cm}
  {\Large\bfseries MA1001B - Modelación Estadística para la Toma de Decisiones\par}
  \vspace{0.7cm}
  {\large Lección 40: Bondad de ajuste ji-cuadrada\par}
  \vspace{0.7cm}
  {\color{TecRojo}\rule{\textwidth}{1.2pt}}
  \vspace{0.8cm}

  Facultad MA1001B\\[0.3cm]
  Periodo\\[0.3cm]
  Tecnol\'ogico de Monterrey
\end{frame}

\begin{frame}{La pregunta de la mezcla afirmada}
  \begin{alertblock}{Pregunta guía}
    ¿Cómo contrastamos si $n=200$ tickets siguen la mezcla hipotética
    $0.40$, $0.30$, $0.20$, $0.10$?
  \end{alertblock}

  \vspace{0.6em}
  Al final de esta lección, usted puede:
  \begin{itemize}
    \item formar conteos esperados $E=np$;
    \item calcular $\chi^2=\sum(O-E)^2/E$ con $df=k-1$;
    \item leer un p-valor de cola derecha a $\alpha=0.05$;
    \item explicar por qué $E<5$ hace pobre la aproximación.
  \end{itemize}

  \vfill
  \scriptsize Todos los conteos de tickets usados hoy son sintéticos. No se usan datos reales de empresa.
\end{frame}

\begin{frame}{Una mezcla hipotética de cuatro categorías}
  \small
  Un mostrador de soporte sintético afirma que los tickets llegan como 40
  por ciento teléfono, 30 por ciento chat, 20 por ciento correo y 10 por
  ciento aplicación. La muestra es $n=200$ con conteos observados $100$,
  $50$, $30$ y $20$.

  \[
    E_i = np_i \quad\Rightarrow\quad 80,\ 60,\ 40,\ 20.
  \]

  \begin{exampleblock}{Hipótesis de bondad de ajuste}
    $H_0$: la mezcla iguala $(0.40,0.30,0.20,0.10)$ frente a $H_1$: la
    mezcla difiere, con $\alpha=0.05$ y $df=3$.
  \end{exampleblock}
\end{frame}

\begin{frame}[fragile]{Paso 1: Observados frente a esperados}
  \begin{columns}[T]
    \begin{column}{0.42\textwidth}
      \small
      \begin{lstlisting}
p = [0.40, 0.30, 0.20, 0.10]
expected = n * p
      \end{lstlisting}

      \begin{block}{Salida verificada}
        Observados: $100$, $50$, $30$, $20$\\
        Esperados: $80$, $60$, $40$, $20$\\
        Todos $E\ge 5$: sí
      \end{block}
    \end{column}
    \begin{column}{0.58\textwidth}
      \centering
      \includegraphics[width=\textwidth]{../figuras/bondad_ajuste_01_observados_esperados.png}
      \vspace{0.2em}
      \scriptsize Conteos observados frente a esperados del Paso 1
    \end{column}
  \end{columns}
\end{frame}

\begin{frame}{Ji-cuadrada con $df=3$}
  \small
  \[
    \chi^2=\frac{(100-80)^2}{80}+\frac{(50-60)^2}{60}
    +\frac{(30-40)^2}{40}+\frac{(20-20)^2}{20}=9.166667.
  \]
  \[
    df=4-1=3,\qquad p=P(\chi^2_3\ge 9.166667)=0.027155.
  \]
  Como $\chi^2>\chi^{2*}=7.814728$ y $p<\alpha$, la prueba rechaza $H_0$.
\end{frame}

\begin{frame}[fragile]{Paso 2: La decisión de bondad de ajuste}
  \begin{columns}[T]
    \begin{column}{0.40\textwidth}
      \small
      \begin{lstlisting}
contrib = (O - E)**2 / E
chi2 = contrib.sum()
p = chi2_dist.sf(chi2, 3)
      \end{lstlisting}

      \begin{block}{Salida verificada}
        $\chi^2=9.166667$, $df=3$\\
        $p=0.027155$\\
        Decisión: rechazar $H_0$
      \end{block}
    \end{column}
    \begin{column}{0.60\textwidth}
      \centering
      \includegraphics[width=\textwidth]{../figuras/bondad_ajuste_02_chi_cuadrada.png}
      \vspace{0.2em}
      \scriptsize Contribuciones celulares del Paso 2
    \end{column}
  \end{columns}
\end{frame}

\begin{frame}{Paso 3: $E<5$ hace pobre la aproximación}
  \begin{columns}[T]
    \begin{column}{0.42\textwidth}
      \small
      Un seguimiento apresurado $n=20$ tiene conteos esperados $8$, $6$,
      $4$ y $2$.

      \vspace{0.3em}
      Dos celdas tienen $E<5$. El conteo observado de App es $0$.

      \vspace{0.3em}
      Esa celda App aporta $2.000000$ si $O=0$ y $0.500000$ si $O=1$.

      \vspace{0.5em}
      \textbf{Límite:} un ticket mueve $\chi^2$ en $1.5$ cuando $E=2$. La
      curva de muestra grande no es aquí una herramienta estable.
    \end{column}
    \begin{column}{0.58\textwidth}
      \centering
      \includegraphics[width=\textwidth]{../figuras/bondad_ajuste_03_limite_esperados_pequenos.png}
      \vspace{0.2em}
      \scriptsize Conteos esperados menores que 5 del Paso 3
    \end{column}
  \end{columns}
\end{frame}

\begin{frame}{Verificación de conceptos}
  \begin{enumerate}
    \item Calcule el conteo esperado de teléfono a partir de $n=200$ y
      $p=0.40$.
    \item ¿Por qué la bondad de ajuste es de cola derecha?
    \item Si $\chi^2=9.166667$, $df=3$ y $p=0.027155$, ¿cuál es la
      decisión a $\alpha=0.05$?
    \item ¿Por qué $E=2$ es un problema aunque el software imprima un
      p-valor?
  \end{enumerate}

  \vspace{0.6em}
  \begin{block}{Conexión con la Lección 41}
    La sesión puede continuar directamente con una prueba de independencia
    ji-cuadrada en una tabla de contingencia de doble entrada.
  \end{block}

  \vfill
  \scriptsize Fuente: OpenStax, \textit{Introductory Business Statistics 2e},
  Capítulo 11, Sección 11.3. CC BY-NC-SA.
\end{frame}

\appendix



\end{document}
